\begin{longtable}{p{1.7cm} p{2.3cm} p{1.9cm} p{5.0cm} p{1.1cm} p{1.9cm}}
\caption{Requerimientos funcionales de Bases (RT)}\label{tab:bt_rf}\\
\toprule
\textbf{ID} & \textbf{T\'itulo} & \textbf{Actor} & \textbf{Descripci\'on} & \textbf{Prioridad} & \textbf{Origen} \\
\midrule
\endfirsthead
\toprule
\textbf{ID} & \textbf{T\'itulo} & \textbf{Actor} & \textbf{Descripci\'on} & \textbf{Prioridad} & \textbf{Origen} \\
\midrule
\endhead
\bottomrule
\endlastfoot
\textbf{RF-13.01} & Declaración de arquitectura multicapa & Equipo de proyecto / Arquitecto & El PROPONENTE debe declarar la arquitectura lógica de la solución organizada en las 8 capas del modelo de referencia: Presentación, Borde y Exposición, Puerta de Enlace, Servicios de Negocio, Integración y Eventos, Datos, Seguridad Transversal y Observabilidad Transversal. & Muy Alta & RT-02.01, RT-02.03 \\
\textbf{RF-13.02} & Tabla de emplazamiento nube/on-premise & Equipo de proyecto & El PROPONENTE debe presentar una tabla de emplazamiento que asigne cada componente de la solución a nube o a on-premise, justificando cada decisión según el criterio de emplazamiento (latencia, criticidad, volumen de datos, acoplamiento físico, elasticidad, regulación). & Muy Alta & Cap. 3.1, RT-03.01 \\
\textbf{RF-13.03} & Registro de decisiones de arquitectura (ADR) & Equipo de proyecto & El PROPONENTE debe mantener un registro de decisiones de arquitectura (ADR) fechado, con la alternativa escogida, las alternativas descartadas y el criterio de decisión, actualizado durante toda la ejecución. & Alta & RT-02.04 \\
\textbf{RF-13.04} & Modelo de dominio de negocio & Equipo de proyecto & El PROPONENTE debe presentar un modelo de dominio del negocio del caso, con las entidades principales, sus relaciones y los eventos de negocio que las modifican. & Muy Alta & RT-02.13 \\
\textbf{RF-14.01} & Clasificación de información por nivel de sensibilidad & Administrador de seguridad & El sistema debe permitir clasificar la información del CLIENTE por nivel de sensibilidad (pública, interna, confidencial, restringida), con controles diferenciados por nivel documentados en una matriz. & Muy Alta & RT-11.03 \\
\textbf{RF-14.02} & Matriz de controles de seguridad trazable a ISO 27001/27002 & Administrador de seguridad & El PROPONENTE debe mantener una matriz de controles de seguridad trazable a ISO/IEC 27001 e ISO/IEC 27002, indicando el control, su implementación concreta en la solución y la evidencia que lo acredita. & Muy Alta & RT-11.05 \\
\textbf{RF-14.03} & Declaración de superficie de exposición & Administrador de seguridad & El PROPONENTE debe declarar la superficie de exposición completa de la solución: cada nombre de dominio, puerto y servicio alcanzable desde fuera de la red del CLIENTE. & Alta & RT-11.13 \\
\textbf{RF-14.04} & Plan de respuesta a incidentes de seguridad & Administrador de seguridad & El PROPONENTE debe disponer de un plan de respuesta a incidentes de seguridad que defina clasificación, cadena de escalamiento, plazos, responsables y protocolo de comunicación al CLIENTE dentro de las 2 horas de detectado un incidente de severidad crítica. & Muy Alta & RT-11.18 \\
\textbf{RF-14.05} & Notificación de brecha de seguridad al CLIENTE & Administrador de seguridad & El sistema (o el equipo de seguridad del ADJUDICATARIO) debe notificar al CLIENTE cualquier brecha de seguridad o de datos personales dentro de las 24 horas de su detección, con informe preliminar, y entregar el análisis de causa raíz dentro de los 5 días hábiles siguientes. & Muy Alta & RT-11.19 \\
\textbf{RF-15.01} & Gestión centralizada de identidad & Administrador del sistema & El sistema debe contar con una gestión de identidad centralizada, con federación mediante OpenID Connect y OAuth 2.1, o SAML 2.0 cuando la integración con el CLIENTE lo requiera, e integración con el directorio corporativo del CLIENTE por LDAP o equivalente en la nube. & Muy Alta & RT-12.01 \\
\textbf{RF-15.02} & Inicio de sesión único (SSO) para todos los módulos & Administrador del sistema & El sistema debe implementar inicio de sesión único para todos los módulos de la solución, con cierre de sesión propagado a todos ellos. & Muy Alta & RT-12.02 \\
\textbf{RF-15.03} & Autenticación multifactor (MFA) para administradores y acceso externo & Administrador del sistema & El sistema debe exigir autenticación multifactor para personas usuarias administradoras, para todo acceso privilegiado y para todo acceso originado fuera de la red corporativa, soportando factores FIDO2 o claves de acceso. & Muy Alta & RT-12.03 \\
\textbf{RF-15.04} & Control de acceso basado en roles (RBAC) y segregación de funciones & Administrador del sistema & El sistema debe implementar control de acceso basado en roles, complementado con control basado en atributos donde el proceso lo exija, con matriz de segregación de funciones documentada y verificable. & Muy Alta & RT-12.05 \\
\textbf{RF-15.05} & Aprovisionamiento y desaprovisionamiento automatizado de cuentas & Administrador del sistema & El sistema debe automatizar el aprovisionamiento y desaprovisionamiento de cuentas, ligado al ciclo de vida laboral, con baja efectiva en un plazo no superior a 24 horas desde la desvinculación. & Alta & RT-12.10 \\
\textbf{RF-16.01} & Observabilidad unificada para nube y on-premise & Administrador del sistema / Equipo de operaciones & El sistema debe implementar observabilidad unificada para nube y on-premise, con métricas, registros y trazas distribuidas correlacionadas por un identificador único de transacción, instrumentadas conforme a OpenTelemetry. & Muy Alta & RT-14.01 \\
\textbf{RF-16.02} & Acceso del CLIENTE a tableros operacionales y de negocio & Administrador del sistema / CLIENTE & El CLIENTE debe disponer de acceso propio y permanente a los tableros operacionales y de negocio, con datos en tiempo real y capacidad de exportación. & Muy Alta & RT-14.02 \\
\textbf{RF-16.03} & Alertamiento basado en síntomas de negocio & Administrador del sistema / Equipo de operaciones & El sistema debe implementar alertamiento basado en síntomas de negocio y no sólo en umbrales de infraestructura, con supresión de ruido, agrupación, escalamiento automático y turnos de disponibilidad declarados. & Alta & RT-14.04 \\
\textbf{RF-16.04} & Análisis de causa raíz obligatorio tras incidente crítico & Equipo de operaciones / Gerente de servicio & El sistema (o el equipo del ADJUDICATARIO) debe ejecutar un análisis de causa raíz obligatorio tras cada incidente crítico, con informe entregable dentro de cinco días hábiles y seguimiento de las acciones correctivas hasta su cierre. & Alta & RT-14.06 \\
\textbf{RF-17.01} & Módulo de administración de usuarios, roles y permisos & Administrador del sistema & El sistema debe disponer de un módulo de administración que permita al CLIENTE, sin intervención del ADJUDICATARIO, gestionar personas usuarias, roles, permisos, unidades organizacionales y sus jerarquías. & Muy Alta & RT-16.01 \\
\textbf{RF-17.02} & Parametrización de reglas de negocio desde interfaz de administración & Administrador del sistema & El sistema debe permitir parametrizar reglas de negocio (umbrales, plazos, montos, tolerancias, catálogos, listas de valores, textos de notificación) desde la interfaz de administración, con control de versiones y registro de quién cambió qué y cuándo. & Muy Alta & RT-16.02 \\
\textbf{RF-17.03} & Aprobación de dos perfiles para cambios de parámetros con impacto operacional & Administrador del sistema & El sistema debe requerir aprobación de un segundo perfil para todo cambio de parámetro con impacto operacional, quedando registrado con su justificación. & Alta & RT-16.03 \\
\textbf{RF-17.04} & Registro de auditoría inalterable & Administrador del sistema & El sistema debe registrar toda operación que cree, modifique o elimine información, con identificación de la persona o sistema que la ejecutó, fecha/hora con zona horaria, origen, valores anteriores y posteriores. El registro de auditoría debe ser inalterable y no modificable por ningún perfil. & Muy Alta & RT-16.06, RT-16.07 \\
\textbf{RF-17.05} & Consulta y exportación de auditoría desde la interfaz & Administrador del sistema / Auditor & El CLIENTE debe poder consultar y exportar la auditoría desde la interfaz, con filtros por persona, período, entidad y tipo de operación, sin requerir acceso a la base de datos. & Muy Alta & RT-16.08 \\
\textbf{RF-17.06} & Soporte de flujos de trabajo configurables & Administrador del sistema & El sistema debe soportar flujos de trabajo con estados, transiciones, responsables, plazos, escalamiento automático por vencimiento y delegación por ausencia, configurables por el CLIENTE sin desarrollo (al menos en responsables, plazos y niveles de aprobación). & Alta & RT-16.11, RT-16.12 \\
\textbf{RF-17.07} & Bandeja de tareas unificada & Persona usuaria (interna) & El sistema debe proveer una bandeja de tareas unificada para cada persona usuaria, donde se visualicen todas sus tareas pendientes (aprobaciones, revisiones, excepciones) con priorización y alerta de vencimiento. & Alta & RT-16.13 \\
\textbf{RF-17.08} & Gestión documental con versionado y control de acceso & Administrador del sistema & El sistema debe gestionar documentos con versionado, metadatos, control de acceso, búsqueda por contenido y por metadato, y previsualización sin descarga. Los documentos deben almacenarse cifrados con verificación de integridad. & Alta & RT-16.15, RT-16.16 \\
\textbf{RF-17.09} & Generación de documentos a partir de plantillas & Administrador del sistema & El sistema debe generar documentos a partir de plantillas administrables por el CLIENTE, con datos de la transacción y salida en formato abierto PDF, DOCX, EXCELL. & Media & RT-16.19 \\
\textbf{RF-17.10} & Notificaciones multicanal (correo, SMS, notificación en app) & Administrador del sistema & El sistema debe enviar notificaciones por al menos tres canales: correo electrónico, notificación en la aplicación y SMS o mensajería instantánea, con plantillas administrables por el CLIENTE. & Alta & RT-16.20, RT-16.21 \\
\textbf{RF-17.11} & Búsqueda global con indexación de texto completo & Persona usuaria (interna) & El sistema debe disponer de búsqueda global con indexación de texto completo con tolerancia a errores de escritura, filtros facetados y respeto del control de acceso de la persona que busca. & Alta & RT-16.27 \\
\textbf{RF-17.12} & Listados ordenables, filtrables, paginados y exportables & Persona usuaria (interna) & El sistema debe presentar listados ordenables, filtrables, paginados y exportables en formatos abiertos, con el filtro aplicado reflejado en la exportación. & Alta & RT-16.28 \\
\textbf{RF-17.13} & Portal público con catálogo y contacto (sin autenticación) & Usuario no autenticado & El sistema debe disponer de un portal público con la información que el caso determine (catálogo sin precios/stock, contacto, información institucional), accesible sin autenticación, con estándares de accesibilidad y desempeño. & Media & RT-16.31 \\
\textbf{RF-18.01} & Soporte de firma electrónica avanzada (Ley 19.799) & Administrador del sistema / Persona usuaria & El sistema debe soportar firma electrónica conforme a la Ley N\ensuremath{^\circ} 19.799, con firma avanzada para los actos que el caso lo requiera (guía de despacho electrónica, acuse de recibo, contratos), y verificación de validez del certificado al momento de la firma, generando sello de tiempo y conservando la evidencia de firma. & Muy Alta & RT-16.17, RT-16.18 \\
\textbf{RF-18.02} & Configuración de preferencias de notificación por usuario & Persona usuaria (interna o externa) & Cada persona usuaria debe poder configurar sus preferencias de canal (correo, SMS, notificación en app) y de frecuencia para las notificaciones, respetando las que el CLIENTE defina como obligatorias. & Media & RT-16.22 \\
\textbf{RF-18.03} & Envío asíncrono de notificaciones con reintento y registro & Administrador del sistema & El sistema debe enviar notificaciones de forma asíncrona, con reintento ante falla, control de duplicados y registro de entrega, apertura y error por cada mensaje. & Alta & RT-16.23 \\
\textbf{RF-13.04} & Modelo de dominio de negocio & Equipo de proyecto & El PROPONENTE debe presentar un modelo de dominio del negocio del caso, con las entidades principales, sus relaciones y los eventos de negocio que las modifican. & Uso de la expresión ambigua "o similar" al referirse a la notación del modelo de dominio (viola la regla de evitar frases que dejan abierta la interpretación). No es verificable qué otra notación sería aceptada. Se recomienda: "...entregado con notación UML". &  \\
\textbf{RF-17.09} & Generación de documentos a partir de plantillas & Administrador del sistema & El sistema debe generar documentos a partir de plantillas administrables por el CLIENTE, con datos de la transacción y salida en formato abierto (PDF, DOCX, etc.). & Contiene "etc." en "formato abierto (PDF, DOCX, etc.)", lo que viola explícitamente la regla de no usar "etc." en la redacción de requerimientos (deja el alcance indefinido). Se recomienda enumerar taxativamente los formatos aceptados. &  \\
\textbf{RF-17.11} & Búsqueda global con indexación de texto completo & Persona usuaria (interna) & El sistema debe disponer de búsqueda global con indexación de texto completo, tolerancia a errores de escritura, filtros facetados y respeto del control de acceso de la persona que busca. & El resultado esperado usa el adjetivo ambiguo "búsqueda rápida" sin métrica asociada, lo que lo hace no verificable (viola la regla de evitar adjetivos ambiguos tipo "rápido"). Se recomienda fijar un umbral, ej. "tiempo de respuesta $\le$ 2 segundos para el percentil 95". &  \\
\end{longtable}

